\section{Lectura de artículo.}
\subsection*{a. Resumen.}
\addcontentsline{toc}{subsection}{a. Resumen}

Hoy en día, gran parte de nuestra información personal es recolectada por dispositivos conectados a la red, dejándola sujeta a la oferta y demanda. Si bien esto puede influir en la confianza de los usuarios, el tratar de medir la confianza en las empresas sin conocer el valor de nuestra información en el mercado puede mostrarnos una percepción errónea de nuestro rol en la comunidad virtual.

En la sociedad actual, siendo guiada por el manejo datos, las compañías dedicadas a la tecnología lideran el mercado digital por la forma de gestionar la información con la que se miden transacciones diarias de dinero, ideas, bienes, etc., las cuales, en conjunto y día tras día, son capaces de proyectar gran parte de nuestra vida.

Nuestra existencia en el Internet se ve confirmada por cada clic que damos, lo cual refleja comportamientos cuyos patrones e incluso variaciones son medibles con la cantidad de información correcta. Con esto, es posible calcular las probabilidades de comportamiento de los usuarios a futuro y, con ello, mostrar anuncios más acordes a sus intereses.

Es por esto que, en aplicaciones que ofrecen servicios aparentemente gratuitos (como las redes sociales), los ingresos de las empresas provienen de vender los datos de sus usuarios a terceros y mostrar mensajes patrocinados, fundamentando la riqueza de estas empresas en la cantidad de usuarios conectados $-$y de los datos que a su vez extraen$-$ y dando pie al desarrollo de una competencia entre entidades por acaparar, privatizar y explotar aquello que empresarialmente se ha denominado como “el nuevo petróleo”.

Además puesto que, de la inmensa cantidad de datos que generan los usuarios, solo la porción disponible para su estudio es la que se procesa en información aprovechable, la industria de la información crece a la par de los avances en tecnologías de procesamiento de datos cada vez más eficientes para aprovechar los conjuntos de Big Data al máximo posible.

Sin embargo, toda esta industria de la mercantilización de la información de los consumidores se da a costa de una constante reducción de su privacidad, puesto se busca recopilar la mayor cantidad de datos posible de éstos, siendo ésta una tarea cada vez más exhaustiva dada la creciente presencia del Internet de las Cosas (IoT) en la cotidianeidad, como en electrodomésticos, joyería y asistentes de IA con sensores integrados.

El IoT se ha vuelto tan importante que, pese a que la mayor parte de la información recolectada y procesada es de carácter personal, a mediados de la década del 2010, la Unión Internacional de Comunicaciones lo consideró la base de la red global de información, y la Fundación Ellen MacArthur predecía que para 2022 la red del IoT se conformaría por hasta 1000 billones de dispositivos, todos con la funcionalidad de recolección de datos acerca de nosotros, adquiriendo mayor precisión a la par que valor económico.

Según la Organización para la Cooperación y el Desarrollo Económicos, si bien no existe una medida fija que calcule el valor de la información de un usuario, las dos principales perspectivas de evaluación son:
\begin{enumerate}
    \item \textbf{Del mercado:} es posible hacerse la idea del valor de un usuario por los ingresos publicitarios que genera; p.ej. Google lo estimó en 2014 en un promedio de $~$45 dólares.
    \item \textbf{Del usuario:} el valor de la información que define la identidad de uno mismo resulta muy variante, desde menos de una decena de euros hasta algunos miles, lo que evidencia la división de opiniones sobre el valor de la confidencialidad de los datos personales. 
\end{enumerate}

El crecimiento económico debe poder sustentarse en una confianza sólida del usuario hacia el mercado, ideal que se promueve actualmente con la existencia de proyectos orientados al usuario, como ciertas reformas de varios países en la regulación de la protección de datos, y tecnologías como las Tiendas de Información Personal y el Hub of All Things, que permiten intercambios de información personal basados en acuerdos de consentimiento explícito.

Así, observando que la facilidad con la que nuestra información personal diariamente se recolecta puede contrastar con nuestros modelos de privacidad y comodidad, así como las posturas al respecto, concluimos que es importante considerar las consecuencias de mostrar nuestros en las redes digitales y los estándares con los que permitimos su posesión por terceros, esto con el fin de comenzar a tener un mayor control de nuestros datos y del valor que opinamos que tienen.
